\documentclass{report}
\usepackage{amsmath,amssymb}
\setlength{\parindent}{0mm}
\setlength{\parskip}{1em}
\begin{document}
\begin{center}
\rule{6in}{1pt} \
{\large Eric Phipps \\
{\bf Preliminary Investigations of Multi-Core Accelerated Embedded Uncertainty Quantification}}

Sandia National Laboratories \\ PO Box 5800 MS-1318 \\ Albuquerque NM 87185
\\
{\tt etphipp@sandia.gov}\\
H. Carter Edwards\\
Jonathan Hu\end{center}

Quantifying uncertainties in computational simulations is a key challenge
for predictive computational science. An important task within
uncertainty quantification is the propagation of input data uncertainty
to the corresponding simulation output quantities of interest. While many
uncertainty propagation approaches have been investigated throughout the
literature, one approach of interest is the embedded stochastic Galerkin
method. Unlike sampling-based methods that only require evaluating the
forward simulation at an independent set of realizations of the uncertain
input data, and thus decouple the corresponding deterministic and
stochastic discretizations, stochastic Galerkin methods couple these
degrees-of-freedom together. For a given level of accuracy, they trade
fewer stochastic degrees-of-freedom for the additional burden of
formulating and solving the resulting coupled (and much larger) system.
When applied to partial differential equations (PDEs) with random input
data, they have generally been found to efficient for certain types of
linear PDEs, but inefficient for PDEs with nonlinear dependence on either
the random variables or the unknown solution. This difference is
attributed to the often dramatically greater number of non-zeros in the
resulting stochastic Galerkin operator for the nonlinear case.

The stochastic Galerkin operator has a Kronecker product structure
arising from the stochastic and deterministic discretizations, and thus
it is typically organized hierarchically as a block-operator where the
sparsity between the blocks is dictated by the stochastic discretization
and each block has the structure of the deterministic discretization.
This allows for greater code reuse when modifying a code to implement the
stochastic Galerkin method, as well as reuse of solvers and
preconditioners designed for the deterministic problem. However, as
described above, this block operator is much denser in the nonlinear
case, and can even be fully block-dense. Therefore in this work, we
investigate commuting this layout to arrive at a sparse block operator
arising from the deterministic discretization where each (possibly dense)
block has the structure of the stochastic discretization. We investigate
the amenability of this layout to emerging multi-core architectures
containing a high-degree of on-node shared-memory parallelism. We develop
several approaches for implementing the stochastic Galerkin matrix-vector
product using threading for the stochastic blocks, and study their
performance on contemporary multi-core CPUs and many-core GPUs.


\end{document}
